\section{Introduction}

The digitalization of houseplant care represents a growing intersection between the Internet of Things (IoT) and daily domestic life. As urbanization increases and living spaces become denser, the desire to maintain greenery has led to an emerging market for smart gardening devices. A common limitation of existing products is their operation as isolated ecosystems; while they may provide localized alerts for individual plants, they often lack broader community interaction. The Plant Up! ecosystem was developed specifically to address this gap. By combining continuous environmental monitoring with a collaborative, gamified social platform, Plant Up! aims to transition routine plant care from an isolated task into a data-driven, shared social experience.

\subsection{Architectural Challenges: Latency and Dependability}

Designing a social smart gardening platform requires a robust backend infrastructure capable of simultaneous management of continuous sensor data and complex user interactions. The primary technical challenges center on two factors: managing network latency and ensuring system availability.

If an architecture relies exclusively on a centralized cloud server for processing every data point, it inherently introduces Round-Trip Time (RTT) delays. In a traditional configuration, a device records sensor data, transmits the payload to a remote server, and waits for the server to evaluate thresholds before receiving a response. This frequent communication can degrade real-time responsiveness. Furthermore, if network connectivity is lost, the system's functionality is compromised due to a lack of localized autonomous processing.

These challenges form the architectural focus of this thesis, which investigates the trade-offs between Edge-Centric and Cloud-Centric processing paradigms. The objective is to optimize both system responsiveness and long-term reliability. A detailed theoretical analysis of these models is provided in Section 3.2 .

\subsection{Research Objective}

This research investigates the fundamental differences between Edge and Cloud processing within a modular backend architecture, analytically isolating the trade-offs present when implementing large-scale environmental monitoring, as demonstrated by the Plant Up! ecosystem.

To provide a concrete evaluation, Plant Up! serves as a comparative case study. The research isolates the processing logic to analyze how shifting it between the edge device and the cloud impacts latency, infrastructure complexity, and system resilience. The functional domain, houseplant monitoring, provides a specific testing environment. Soil moisture and ambient temperature are metrics that respond slowly to environmental changes. Since this domain does not require the millisecond-level reaction times of industrial machinery, the boundaries of cloud processing can be further explored without compromising the user experience.

\subsection{Scope and Limitations}

\subsubsection{Computational Constraints of the Edge Device}

The physical edge device utilized in this deployment is the ESP32 microcontroller, which is characterized by specific processing limitations and a constrained memory footprint. Due to these hardware constraints, the architecture must carefully distribute tasks between local data aggregation and cloud-based logic. The precise hardware characteristics, including network limitations and energy consumption, are detailed in Chapter 4. This section focuses on how these physical realities influence the placement of processing logic.

\subsubsection{Abstraction of the Transport Layer}

To ensure an isolated comparative analysis, this thesis establishes rigorous scope boundaries. The evaluation focuses specifically on the location of the application logic (Edge versus Cloud computing). To maintain environmental consistency, the transport layer is abstracted, with the Message Queuing Telemetry Transport (MQTT) protocol utilized as the standard for all variants. Alternative application-layer protocols, such as CoAP or HTTP, are excluded from this analysis. By standardizing the communication pathway, the architectural comparison reflects only the impact of the computational placement.
